% Options for packages loaded elsewhere
% Options for packages loaded elsewhere
\PassOptionsToPackage{unicode}{hyperref}
\PassOptionsToPackage{hyphens}{url}
%
\documentclass[
  ignorenonframetext,
  aspectratio=169,
  spanish,
]{beamer}
\newif\ifbibliography
\usepackage{pgfpages}
\setbeamertemplate{caption}[numbered]
\setbeamertemplate{caption label separator}{: }
\setbeamercolor{caption name}{fg=normal text.fg}
\beamertemplatenavigationsymbolshorizontal
% remove section numbering
\setbeamertemplate{part page}{
  \centering
  \begin{beamercolorbox}[sep=16pt,center]{part title}
    \usebeamerfont{part title}\insertpart\par
  \end{beamercolorbox}
}
\setbeamertemplate{section page}{
  \centering
  \begin{beamercolorbox}[sep=12pt,center]{section title}
    \usebeamerfont{section title}\insertsection\par
  \end{beamercolorbox}
}
\setbeamertemplate{subsection page}{
  \centering
  \begin{beamercolorbox}[sep=8pt,center]{subsection title}
    \usebeamerfont{subsection title}\insertsubsection\par
  \end{beamercolorbox}
}
% Prevent slide breaks in the middle of a paragraph
\widowpenalties 1 10000
\raggedbottom
\AtBeginPart{
  \frame{\partpage}
}
\AtBeginSection{
  \ifbibliography
  \else
    \frame{\sectionpage}
  \fi
}
\AtBeginSubsection{
  \frame{\subsectionpage}
}
\usepackage{iftex}
\ifPDFTeX
  \usepackage[T1]{fontenc}
  \usepackage[utf8]{inputenc}
  \usepackage{textcomp} % provide euro and other symbols
\else % if luatex or xetex
  \usepackage{unicode-math} % this also loads fontspec
  \defaultfontfeatures{Scale=MatchLowercase}
  \defaultfontfeatures[\rmfamily]{Ligatures=TeX,Scale=1}
\fi
\usepackage{lmodern}

\usetheme[]{default}
\usefonttheme[]{default}
\ifPDFTeX\else
  % xetex/luatex font selection
\fi
% Use upquote if available, for straight quotes in verbatim environments
\IfFileExists{upquote.sty}{\usepackage{upquote}}{}
\IfFileExists{microtype.sty}{% use microtype if available
  \usepackage[]{microtype}
  \UseMicrotypeSet[protrusion]{basicmath} % disable protrusion for tt fonts
}{}
\makeatletter
\@ifundefined{KOMAClassName}{% if non-KOMA class
  \IfFileExists{parskip.sty}{%
    \usepackage{parskip}
  }{% else
    \setlength{\parindent}{0pt}
    \setlength{\parskip}{6pt plus 2pt minus 1pt}}
}{% if KOMA class
  \KOMAoptions{parskip=half}}
\makeatother

\usepackage{color}
\usepackage{fancyvrb}
\newcommand{\VerbBar}{|}
\newcommand{\VERB}{\Verb[commandchars=\\\{\}]}
\DefineVerbatimEnvironment{Highlighting}{Verbatim}{commandchars=\\\{\}}
% Add ',fontsize=\small' for more characters per line
\usepackage{framed}
\definecolor{shadecolor}{RGB}{241,243,245}
\newenvironment{Shaded}{\begin{snugshade}}{\end{snugshade}}
\newcommand{\AlertTok}[1]{\textcolor[rgb]{0.68,0.00,0.00}{#1}}
\newcommand{\AnnotationTok}[1]{\textcolor[rgb]{0.37,0.37,0.37}{#1}}
\newcommand{\AttributeTok}[1]{\textcolor[rgb]{0.40,0.45,0.13}{#1}}
\newcommand{\BaseNTok}[1]{\textcolor[rgb]{0.68,0.00,0.00}{#1}}
\newcommand{\BuiltInTok}[1]{\textcolor[rgb]{0.00,0.23,0.31}{#1}}
\newcommand{\CharTok}[1]{\textcolor[rgb]{0.13,0.47,0.30}{#1}}
\newcommand{\CommentTok}[1]{\textcolor[rgb]{0.37,0.37,0.37}{#1}}
\newcommand{\CommentVarTok}[1]{\textcolor[rgb]{0.37,0.37,0.37}{\textit{#1}}}
\newcommand{\ConstantTok}[1]{\textcolor[rgb]{0.56,0.35,0.01}{#1}}
\newcommand{\ControlFlowTok}[1]{\textcolor[rgb]{0.00,0.23,0.31}{\textbf{#1}}}
\newcommand{\DataTypeTok}[1]{\textcolor[rgb]{0.68,0.00,0.00}{#1}}
\newcommand{\DecValTok}[1]{\textcolor[rgb]{0.68,0.00,0.00}{#1}}
\newcommand{\DocumentationTok}[1]{\textcolor[rgb]{0.37,0.37,0.37}{\textit{#1}}}
\newcommand{\ErrorTok}[1]{\textcolor[rgb]{0.68,0.00,0.00}{#1}}
\newcommand{\ExtensionTok}[1]{\textcolor[rgb]{0.00,0.23,0.31}{#1}}
\newcommand{\FloatTok}[1]{\textcolor[rgb]{0.68,0.00,0.00}{#1}}
\newcommand{\FunctionTok}[1]{\textcolor[rgb]{0.28,0.35,0.67}{#1}}
\newcommand{\ImportTok}[1]{\textcolor[rgb]{0.00,0.46,0.62}{#1}}
\newcommand{\InformationTok}[1]{\textcolor[rgb]{0.37,0.37,0.37}{#1}}
\newcommand{\KeywordTok}[1]{\textcolor[rgb]{0.00,0.23,0.31}{\textbf{#1}}}
\newcommand{\NormalTok}[1]{\textcolor[rgb]{0.00,0.23,0.31}{#1}}
\newcommand{\OperatorTok}[1]{\textcolor[rgb]{0.37,0.37,0.37}{#1}}
\newcommand{\OtherTok}[1]{\textcolor[rgb]{0.00,0.23,0.31}{#1}}
\newcommand{\PreprocessorTok}[1]{\textcolor[rgb]{0.68,0.00,0.00}{#1}}
\newcommand{\RegionMarkerTok}[1]{\textcolor[rgb]{0.00,0.23,0.31}{#1}}
\newcommand{\SpecialCharTok}[1]{\textcolor[rgb]{0.37,0.37,0.37}{#1}}
\newcommand{\SpecialStringTok}[1]{\textcolor[rgb]{0.13,0.47,0.30}{#1}}
\newcommand{\StringTok}[1]{\textcolor[rgb]{0.13,0.47,0.30}{#1}}
\newcommand{\VariableTok}[1]{\textcolor[rgb]{0.07,0.07,0.07}{#1}}
\newcommand{\VerbatimStringTok}[1]{\textcolor[rgb]{0.13,0.47,0.30}{#1}}
\newcommand{\WarningTok}[1]{\textcolor[rgb]{0.37,0.37,0.37}{\textit{#1}}}

\usepackage{longtable,booktabs,array}
\usepackage{calc} % for calculating minipage widths
\usepackage{caption}
% Make caption package work with longtable
\makeatletter
\def\fnum@table{\tablename~\thetable}
\makeatother
\usepackage{graphicx}
\makeatletter
\newsavebox\pandoc@box
\newcommand*\pandocbounded[1]{% scales image to fit in text height/width
  \sbox\pandoc@box{#1}%
  \Gscale@div\@tempa{\textheight}{\dimexpr\ht\pandoc@box+\dp\pandoc@box\relax}%
  \Gscale@div\@tempb{\linewidth}{\wd\pandoc@box}%
  \ifdim\@tempb\p@<\@tempa\p@\let\@tempa\@tempb\fi% select the smaller of both
  \ifdim\@tempa\p@<\p@\scalebox{\@tempa}{\usebox\pandoc@box}%
  \else\usebox{\pandoc@box}%
  \fi%
}
% Set default figure placement to htbp
\def\fps@figure{htbp}
\makeatother



\ifLuaTeX
\usepackage[bidi=basic]{babel}
\else
\usepackage[bidi=default]{babel}
\fi
% get rid of language-specific shorthands (see #6817):
\let\LanguageShortHands\languageshorthands
\def\languageshorthands#1{}


\setlength{\emergencystretch}{3em} % prevent overfull lines

\providecommand{\tightlist}{%
  \setlength{\itemsep}{0pt}\setlength{\parskip}{0pt}}



 


\setbeamertemplate{headline}{} % Elimina la cabecera del tema Madrid
\setbeamertemplate{section in head/foot}{} % Limpia la numeración de secciones
\renewcommand{\contentsname}{Índice}
\renewcommand{\refname}{Referencias}
\renewcommand{\figurename}{Figura}
\renewcommand{\tablename}{Tabla}
\logo{%
  \begin{tabular}{c}
    \includegraphics[height=1.2cm]{EscudoU.png}
  \end{tabular}%
}
% --- CONFIGURACIÓN DE COLORES INSTITUCIONALES ---
\definecolor{UnicorGreen}{HTML}{024426}
\definecolor{UnicorText}{HTML}{F9F9F9}
\definecolor{UnicorWhite}{HTML}{FFFFFF}
\definecolor{UnicorGold}{HTML}{D1B061}

% Aplicación al fondo y texto general
\setbeamercolor{background canvas}{bg=UnicorGreen}
\setbeamercolor{normal text}{fg=UnicorText}

% Aplicación a títulos y cabeceras
%\setbeamercolor{title}{fg=UnicorWhite}
\setbeamercolor{frametitle}{fg=UnicorWhite}
\setbeamercolor{framesubtitle}{fg=UnicorWhite}

% Aplicación a viñetas y estructura
\setbeamercolor{structure}{fg=UnicorWhite}
\setbeamercolor{itemize item}{fg=UnicorWhite}
\setbeamercolor{itemize subitem}{fg=UnicorWhite}
\setbeamercolor{enumerate item}{fg=UnicorWhite}

% Forzar a LaTeX a repintar el texto general al iniciar el documento
\AtBeginDocument{\usebeamercolor[fg]{normal text}}

% Aplicación a hipervínculos
\hypersetup{colorlinks=true, urlcolor=UnicorGold, linkcolor=UnicorWhite}

% --- ACTIVAR ESTILO DE BLOQUES Y SOMBRAS (ESTILO MADRID) ---
\useinnertheme[shadow=true]{rounded}

% --- COLORES DEL RECUADRO DEL TÍTULO PRINCIPAL ---
% Fondo dorado con texto verde
\setbeamercolor{title}{bg=UnicorGold, fg=UnicorGreen}
\setbeamercolor{frametitle}{fg=UnicorWhite} % Mantiene el título de las diapositivas normal
\setbeamercolor{framesubtitle}{fg=UnicorWhite}

% --- COLORES PARA TEOREMAS, DEFINICIONES Y BLOQUES ---
% Título del bloque (Fondo dorado, texto verde)
\setbeamercolor{block title}{bg=UnicorGold, fg=UnicorGreen}
% Cuerpo del bloque (Fondo blanco, texto verde)
\setbeamercolor{block body}{bg=UnicorWhite, fg=UnicorGreen}

\AtBeginDocument{
  \author{%
    \begin{tabular}{c}
      Mario A. Morales Rivera \\
      \footnotesize \texttt{mamorales@correo.unicordoba.edu.co} \\
      \footnotesize Universidad de Córdoba
    \end{tabular}
  }
  \institute{}
}
% --- LIMPIAR DIAPOSITIVAS DE SECCIÓN (ELIMINAR "SECTION X") ---
\setbeamertemplate{section page}{
  \begin{center}
    \begin{beamercolorbox}[sep=12pt,center,rounded=true,shadow=true]{title}
      \usebeamerfont{title}\insertsection\par
    \end{beamercolorbox}
  \end{center}
}
\makeatletter
\@ifpackageloaded{tcolorbox}{}{\usepackage[skins,breakable]{tcolorbox}}
\@ifpackageloaded{fontawesome5}{}{\usepackage{fontawesome5}}
\definecolor{quarto-callout-color}{HTML}{909090}
\definecolor{quarto-callout-note-color}{HTML}{0758E5}
\definecolor{quarto-callout-important-color}{HTML}{CC1914}
\definecolor{quarto-callout-warning-color}{HTML}{EB9113}
\definecolor{quarto-callout-tip-color}{HTML}{00A047}
\definecolor{quarto-callout-caution-color}{HTML}{FC5300}
\definecolor{quarto-callout-color-frame}{HTML}{acacac}
\definecolor{quarto-callout-note-color-frame}{HTML}{4582ec}
\definecolor{quarto-callout-important-color-frame}{HTML}{d9534f}
\definecolor{quarto-callout-warning-color-frame}{HTML}{f0ad4e}
\definecolor{quarto-callout-tip-color-frame}{HTML}{02b875}
\definecolor{quarto-callout-caution-color-frame}{HTML}{fd7e14}
\makeatother
\makeatletter
\@ifpackageloaded{caption}{}{\usepackage{caption}}
\AtBeginDocument{%
\ifdefined\contentsname
  \renewcommand*\contentsname{Tabla de contenidos}
\else
  \newcommand\contentsname{Tabla de contenidos}
\fi
\ifdefined\listfigurename
  \renewcommand*\listfigurename{Listado de Figuras}
\else
  \newcommand\listfigurename{Listado de Figuras}
\fi
\ifdefined\listtablename
  \renewcommand*\listtablename{Listado de Tablas}
\else
  \newcommand\listtablename{Listado de Tablas}
\fi
\ifdefined\figurename
  \renewcommand*\figurename{Figura}
\else
  \newcommand\figurename{Figura}
\fi
\ifdefined\tablename
  \renewcommand*\tablename{Tabla}
\else
  \newcommand\tablename{Tabla}
\fi
}
\@ifpackageloaded{float}{}{\usepackage{float}}
\floatstyle{ruled}
\@ifundefined{c@chapter}{\newfloat{codelisting}{h}{lop}}{\newfloat{codelisting}{h}{lop}[chapter]}
\floatname{codelisting}{Listado}
\newcommand*\listoflistings{\listof{codelisting}{Listado de Listados}}
\makeatother
\makeatletter
\makeatother
\makeatletter
\@ifpackageloaded{caption}{}{\usepackage{caption}}
\@ifpackageloaded{subcaption}{}{\usepackage{subcaption}}
\makeatother
\makeatletter
\@ifpackageloaded{tikz}{}{\usepackage{tikz}}
\makeatother
        \newcommand*\circled[1]{\tikz[baseline=(char.base)]{
          \node[shape=circle,draw,inner sep=1pt] (char) {{\scriptsize#1}};}}  
                  

\usepackage{bookmark}
\IfFileExists{xurl.sty}{\usepackage{xurl}}{} % add URL line breaks if available
\urlstyle{same}
\hypersetup{
  pdftitle={Listas},
  pdfauthor={Mario Morales},
  pdflang={es},
  hidelinks,
  pdfcreator={LaTeX via pandoc}}


\title{Listas}
\subtitle{\emph{Estadística con Python}}
\author{Mario Morales}
\date{}

\begin{document}
\frame{\titlepage}

\renewcommand*\contentsname{Tabla de contenidos}
\begin{frame}[allowframebreaks]
  \frametitle{Tabla de contenidos}
  \setcounter{tocdepth}{3}
  \tableofcontents
\end{frame}
\setcounter{tocdepth}{3}
\tableofcontents
}

\section{Introducción}\label{sec-listas}

\begin{frame}{Introducción}
\index{Listas}
\end{frame}

\begin{frame}{Introducción}
\phantomsection\label{introducciuxf3n}
Las listas sirven para agrupar objetos de diferente tipo, aunque
típicamente se usan para alojar datos de la misma naturaleza (números,
caracteres, lógicos, \ldots), son indexadas y, a diferencia de las
tuplas, son \textbf{mutables}, es decir, se les puede agregar, quitar y
cambiar elementos después de creadas. Las listas soportan todas las
características de las tuplas, pero es posible adicionar, remover y
modificar sus elementos.
\end{frame}

\section{Creación de listas}\label{creaciuxf3n-de-listas}

\begin{frame}[fragile]{Creación de listas}
\phantomsection\label{creaciuxf3n-de-listas-1}
\index{Listas!Creación de}

Para crear una lista colocamos los elementos dentro de corchetes de la
siguiente forma:

\begin{Shaded}
\begin{Highlighting}[]
\NormalTok{l1}\OperatorTok{=}\NormalTok{[}\OperatorTok{{-}}\DecValTok{5}\NormalTok{,}\OperatorTok{{-}}\DecValTok{4}\NormalTok{,}\OperatorTok{{-}}\DecValTok{3}\NormalTok{,}\OperatorTok{{-}}\DecValTok{2}\NormalTok{,}\OperatorTok{{-}}\DecValTok{1}\NormalTok{,}\DecValTok{1}\NormalTok{,}\DecValTok{2}\NormalTok{,}\DecValTok{3}\NormalTok{,}\DecValTok{4}\NormalTok{,}\DecValTok{5}\NormalTok{] }\CommentTok{\# lista de números enteros}
\NormalTok{l2}\OperatorTok{=}\NormalTok{[}\DecValTok{2}\NormalTok{,}\StringTok{"a"}\NormalTok{,}\VariableTok{True}\NormalTok{]    }\CommentTok{\#lista con elementos de distinto tipo}
\BuiltInTok{print}\NormalTok{(l2,}\BuiltInTok{type}\NormalTok{(l2))}
\end{Highlighting}
\end{Shaded}

\begin{verbatim}
[2, 'a', True] <class 'list'>
\end{verbatim}
\end{frame}

\begin{frame}[fragile]{Creación de listas (cont)}
\phantomsection\label{creaciuxf3n-de-listas-cont}
También se pueden crear listas a partir de otros objetos previamente
creados, por ejemplo tuplas, con la función \texttt{list()}
\index{Listas!Función!\texttt{list()}} \index{Función!\texttt{list()}}

\begin{Shaded}
\begin{Highlighting}[]
\NormalTok{t}\OperatorTok{=}\NormalTok{(}\DecValTok{2}\NormalTok{,}\DecValTok{3}\NormalTok{,}\DecValTok{4}\NormalTok{,}\DecValTok{5}\NormalTok{,}\DecValTok{6}\NormalTok{) }
\BuiltInTok{list}\NormalTok{(t)  }
\end{Highlighting}
\end{Shaded}

\begin{verbatim}
[2, 3, 4, 5, 6]
\end{verbatim}

\begin{Shaded}
\begin{Highlighting}[]
\BuiltInTok{list}\NormalTok{(}\BuiltInTok{range}\NormalTok{(}\DecValTok{10}\NormalTok{)) }
\end{Highlighting}
\end{Shaded}

\begin{verbatim}
[0, 1, 2, 3, 4, 5, 6, 7, 8, 9]
\end{verbatim}
\end{frame}

\section{Indexado}\label{indexado}

\begin{frame}[fragile]{Indexado de listas}
\phantomsection\label{sec-IndexadoDeListas}
\index{Listas!Indexado de}

\begin{itemize}
\tightlist
\item
  Cada elemento de una lista está asociado con un \emph{índice} que
  refleja su posición en el arreglo.
\item
  Como en las tuplas, el primer elemento tiene índice 0, el segundo
  índice 1 y así sucesivamente.
\item
  Asociados con la lista \texttt{l1} anterior tenemos 10 índices,
  iniciando en 0 y terminando en 9 .
\item
  Para acceder al elemento con índice 3, es decir, al \textbf{cuarto}
  elemento de la lista, escribimos \texttt{l1{[}3{]}} que corresponde a
  \(-2\)
\end{itemize}
\end{frame}

\begin{frame}[fragile]{Indexado de listas, ejemplos}
\phantomsection\label{indexado-de-listas-ejemplos}
\begin{Shaded}
\begin{Highlighting}[]
\BuiltInTok{print}\NormalTok{(l1)}
\end{Highlighting}
\end{Shaded}

\begin{verbatim}
[-5, -4, -3, -2, -1, 1, 2, 3, 4, 5]
\end{verbatim}

\begin{Shaded}
\begin{Highlighting}[]
\NormalTok{l1[}\DecValTok{3}\NormalTok{] }\CommentTok{\# Cuarto elemento de la lista l1}
\end{Highlighting}
\end{Shaded}

\begin{verbatim}
-2
\end{verbatim}
\end{frame}

\begin{frame}[fragile]{Indexado de listas (cont)}
\phantomsection\label{indexado-de-listas-cont}
Análogo a las tuplas, estudiadas en la clase anterior, con las listas es
posible usar \textbf{índices negativos}, lo cual conduce a indexar por
la derecha. \texttt{l{[}-1{]}} corresponde al último elemento de la
lista, \texttt{l{[}-2{]}} corresponde al anterior a \texttt{l{[}-1{]}} y
así sucesivamente.

\begin{Shaded}
\begin{Highlighting}[]
\NormalTok{l1[}\OperatorTok{{-}}\DecValTok{1}\NormalTok{] }\CommentTok{\# Último elemento de la lista l1}
\NormalTok{l1[}\OperatorTok{{-}}\DecValTok{2}\NormalTok{] }\CommentTok{\# Penúltimo elemento de la lista l1}
\end{Highlighting}
\end{Shaded}
\end{frame}

\section{Edición}\label{ediciuxf3n}

\begin{frame}[fragile]{Inserción, adición y borrado de elementos en una
lista}
\phantomsection\label{inserciuxf3n-adiciuxf3n-y-borrado-de-elementos-en-una-lista}
\index{Listas!Edición!Método\texttt{.insert()}}
\index{Método!\texttt{.insert()}}

A diferencia de las tuplas, las listas se pueden modificar insertando,
adicionando o borrando elementos, veamos como se usa el método
\texttt{insert()} para insertar el número \(-6\) como primer elemento de
la lista \texttt{l1}

\begin{Shaded}
\begin{Highlighting}[]
\NormalTok{l1.insert(}\DecValTok{0}\NormalTok{,}\OperatorTok{{-}}\DecValTok{6}\NormalTok{)}
\BuiltInTok{print}\NormalTok{(l1)}
\end{Highlighting}
\end{Shaded}

\begin{verbatim}
[-6, -5, -4, -3, -2, -1, 1, 2, 3, 4, 5]
\end{verbatim}
\end{frame}

\begin{frame}[fragile]{Inserción (cont)}
\phantomsection\label{inserciuxf3n-cont}
El primer argumento en la sentencia \texttt{l1.insert(0,-6)},
corresponde al índice donde se insertará el número dado en el segundo
argumento (\texttt{-6}). La lista \texttt{l1} ahora tiene un elemento
más, \(-6\), en el índice 0 y el número \(-5\), que estaba en este
índice, pasó al índice 1.
\end{frame}

\begin{frame}[fragile]{Adición de elementos}
\phantomsection\label{adiciuxf3n-de-elementos}
Para agregar un elemento al final de la lista se usa el método
\texttt{append()} como se muestra a
continuación\index{Listas!Edición!Método\texttt{append()}}\index{Método!\texttt{append()}}

\begin{Shaded}
\begin{Highlighting}[]
\NormalTok{l1.append(}\DecValTok{6}\NormalTok{) }
\BuiltInTok{print}\NormalTok{(l1)}
\end{Highlighting}
\end{Shaded}

\begin{verbatim}
[-6, -5, -4, -3, -2, -1, 1, 2, 3, 4, 5, 6]
\end{verbatim}

note que se ha agregado \texttt{6} al final de la lista \texttt{l1}.
\end{frame}

\begin{frame}[fragile]{Adición de elementos (cont)}
\phantomsection\label{adiciuxf3n-de-elementos-cont}
Es posible adicionar más de un elemento al final de una lista
concatenándola con otra lista, usando el operador \texttt{+}, análogo a
como se hace con tuplas.

en este caso la lista \texttt{l2} no se modifica, para conseguirlo
usamos

\begin{Shaded}
\begin{Highlighting}[]
\NormalTok{l2}\OperatorTok{=}\NormalTok{l2}\OperatorTok{+}\NormalTok{[}\DecValTok{4}\NormalTok{,}\DecValTok{5}\NormalTok{,}\DecValTok{6}\NormalTok{]  }
\BuiltInTok{print}\NormalTok{(l2)}
\end{Highlighting}
\end{Shaded}

\begin{verbatim}
[2, 'a', True, 4, 5, 6]
\end{verbatim}
\end{frame}

\begin{frame}[fragile]{Borrado}
\phantomsection\label{borrado}
Para borrar un elemento de una lista se usa \texttt{del\ lista{[}i{]}}
\index{Función!\texttt{del}} donde \texttt{i} corresponde al índice del
elemento que se desea eliminar. A manera de ejemplo, borremos el
elemento \texttt{\textquotesingle{}a\textquotesingle{}} de la lista
\texttt{l2}, (\texttt{i=1})

\begin{Shaded}
\begin{Highlighting}[]
\KeywordTok{del}\NormalTok{ l2[}\DecValTok{1}\NormalTok{]  }
\BuiltInTok{print}\NormalTok{(l2)}
\end{Highlighting}
\end{Shaded}

\begin{verbatim}
[2, True, 4, 5, 6]
\end{verbatim}
\end{frame}

\begin{frame}[fragile]{Borrado (cont)}
\phantomsection\label{borrado-cont}
Mientras que \texttt{del} elimina elementos de una lista por posición,
el método \texttt{remove()} \index{Método!\texttt{remove()}} los elimina
por valor. Para ilustrarlo, seguidamente borramos el elemento \texttt{4}
de la lista \texttt{l2}.

\begin{Shaded}
\begin{Highlighting}[]
\NormalTok{l2.remove(}\DecValTok{4}\NormalTok{)  }
\BuiltInTok{print}\NormalTok{(l2)}
\end{Highlighting}
\end{Shaded}

\begin{verbatim}
[2, True, 5, 6]
\end{verbatim}
\end{frame}

\begin{frame}[fragile]{Borrado (cont)}
\phantomsection\label{borrado-cont-1}
Si en la lista hay más de un elemento con el argumento entregado a
\texttt{remove()}, se elimina el de índice más pequeño, es decir, el
primero que aparezca en la lista, veamos el siguiente ejemplo

\phantomsection\label{annotated-cell-12}%
\begin{Shaded}
\begin{Highlighting}[]
\NormalTok{l3}\OperatorTok{=}\NormalTok{[}\DecValTok{2}\NormalTok{,}\DecValTok{5}\NormalTok{,}\DecValTok{4}\NormalTok{,}\OperatorTok{{-}}\DecValTok{1}\NormalTok{,}\DecValTok{4}\NormalTok{,}\DecValTok{3}\NormalTok{,}\DecValTok{0}\NormalTok{,}\OperatorTok{{-}}\DecValTok{2}\NormalTok{] }\hspace*{\fill}\NormalTok{\circled{1}}
\NormalTok{l3.remove(}\DecValTok{4}\NormalTok{) }\hspace*{\fill}\NormalTok{\circled{2}}
\NormalTok{l3}
\end{Highlighting}
\end{Shaded}

\begin{description}
\tightlist
\item[\circled{1}]
El número 4 está en la lista dos veces
\item[\circled{2}]
Se elimina el primer 4
\end{description}

\begin{verbatim}
[2, 5, -1, 4, 3, 0, -2]
\end{verbatim}
\end{frame}

\begin{frame}[fragile]{Borrado (cont)}
\phantomsection\label{borrado-cont-2}
Otro método para eliminar elementos de listas es \texttt{pop()}
\index{Método!\texttt{pop()}} que recibe un índice, elimina el elemento
que corresponde, pero además regresa el valor eliminado

\begin{Shaded}
\begin{Highlighting}[]
\NormalTok{l3.pop(}\DecValTok{2}\NormalTok{) }\CommentTok{\# regresa {-}1 (índice 2)}
\end{Highlighting}
\end{Shaded}

\begin{verbatim}
-1
\end{verbatim}

\begin{Shaded}
\begin{Highlighting}[]
\BuiltInTok{print}\NormalTok{(l3) }\CommentTok{\# l3 queda modificada, sin {-}1}
\end{Highlighting}
\end{Shaded}

\begin{verbatim}
[2, 5, 4, 3, 0, -2]
\end{verbatim}
\end{frame}

\begin{frame}[fragile]{Borrado (cont)}
\phantomsection\label{borrado-cont-3}
Si no se entrega índice, el método \texttt{pop()} remueve (y regresa) el
último elemento de la lista.

\begin{Shaded}
\begin{Highlighting}[]
\NormalTok{l3.pop()  }\CommentTok{\# regresa {-}2 (el último elemento de l3)}
\end{Highlighting}
\end{Shaded}

\begin{verbatim}
-2
\end{verbatim}

\begin{Shaded}
\begin{Highlighting}[]
\BuiltInTok{print}\NormalTok{(l3) }\CommentTok{\# l3 queda modificada, se remueve {-}2   }
\end{Highlighting}
\end{Shaded}

\begin{verbatim}
[2, 5, 4, 3, 0]
\end{verbatim}
\end{frame}

\begin{frame}[fragile]{Longitud de una lista.}
\phantomsection\label{longitud-de-una-lista.}
\index{Listas!Longitud de una}

Cuando queremos recuperar en número de elementos que contiene una lista,
procedemos de manera similar a como se hace con tuplas, usamos la
función \texttt{len()}

\begin{Shaded}
\begin{Highlighting}[]
\BuiltInTok{len}\NormalTok{(l2)}
\end{Highlighting}
\end{Shaded}

\begin{verbatim}
4
\end{verbatim}

como una aplicación, podemos combinar las funciones \texttt{del} y
\texttt{len()} para borrar el último elemento de una lista

\begin{Shaded}
\begin{Highlighting}[]
\KeywordTok{del}\NormalTok{ l2[}\BuiltInTok{len}\NormalTok{(l2)}\OperatorTok{{-}}\DecValTok{1}\NormalTok{] }
\end{Highlighting}
\end{Shaded}
\end{frame}

\section{Métodos asociados a
listas}\label{muxe9todos-asociados-a-listas}

\begin{frame}{Indexación de listas}
\phantomsection\label{sec-IndexacionDeLista}
\index{Listas!Indexación de}\index{Indexación!de Listas}\index{Indexación}

En esta sección se estudian las distintas formas de acceder, recuperar y
cambiar elementos de una lista.
\end{frame}

\begin{frame}[fragile]{Verificar si un elemento pertenece a una lista}
\phantomsection\label{verificar-si-un-elemento-pertenece-a-una-lista}
\index{Listas!Operador!\texttt{in}}

De manera análoga a como se hizo con las tuplas, se puede verificar si
un elemento pertenece a una lista con el operador \texttt{in}. Si el
elemento dado está en la lista se obtiene \texttt{True}, en caso
contrario, \texttt{False}

\begin{Shaded}
\begin{Highlighting}[]
\BuiltInTok{print}\NormalTok{(}\DecValTok{0} \KeywordTok{in}\NormalTok{ l1, }\DecValTok{3} \KeywordTok{in}\NormalTok{ l1)}
\end{Highlighting}
\end{Shaded}

\begin{verbatim}
False True
\end{verbatim}
\end{frame}

\begin{frame}[fragile]{Índice de un elemento.}
\phantomsection\label{uxedndice-de-un-elemento.}
Con el método \texttt{index()} \index{Método!\texttt{index()}} se
recupera el índice que corresponde a un valor en la lista. Si el valor
está en la lista más de una vez, se regresa el primer índice.

\begin{Shaded}
\begin{Highlighting}[]
\BuiltInTok{print}\NormalTok{(l1)}
\end{Highlighting}
\end{Shaded}

\begin{verbatim}
[-6, -5, -4, -3, -2, -1, 1, 2, 3, 4, 5, 6]
\end{verbatim}

\begin{Shaded}
\begin{Highlighting}[]
\NormalTok{l1.index(}\OperatorTok{{-}}\DecValTok{1}\NormalTok{) }\CommentTok{\# {-}1 tiene índice 5 }
\end{Highlighting}
\end{Shaded}

\begin{verbatim}
5
\end{verbatim}

\begin{Shaded}
\begin{Highlighting}[]
\NormalTok{l3}\OperatorTok{=}\NormalTok{[}\DecValTok{2}\NormalTok{,}\DecValTok{5}\NormalTok{,}\DecValTok{4}\NormalTok{,}\OperatorTok{{-}}\DecValTok{1}\NormalTok{,}\DecValTok{4}\NormalTok{,}\DecValTok{3}\NormalTok{,}\DecValTok{0}\NormalTok{,}\OperatorTok{{-}}\DecValTok{2}\NormalTok{] }\CommentTok{\# en esta lista 4 está dos veces }
\NormalTok{l3.index(}\DecValTok{4}\NormalTok{) }\CommentTok{\# índice del primer 4 }
\end{Highlighting}
\end{Shaded}

\begin{verbatim}
2
\end{verbatim}
\end{frame}

\begin{frame}[fragile]{Indexación de listas (cont)}
\phantomsection\label{indexaciuxf3n-de-listas-cont}
\begin{itemize}
\item
  En caso que el argumento entregado al método \texttt{index()} no esté
  en la lista, ocurre un error.
\item
  Por ejemplo, \texttt{l3.index(0.5)} regresa
  \texttt{ValueError:\ 0.5\ is\ not\ in\ list}, esto es lo que se conoce
  técnicamente como una excepción no controlada (uncaught exception).
\item
  Para evitar que tu código se \texttt{reviente}, podrías indagar si el
  elemento está en la lista, en caso afirmativo usas \texttt{index()},
  en caso negativo, imprimes un mensaje
\end{itemize}
\end{frame}

\begin{frame}[fragile]{Indexación de listas (cont)}
\phantomsection\label{indexaciuxf3n-de-listas-cont-1}
\begin{Shaded}
\begin{Highlighting}[]
\NormalTok{i}\OperatorTok{=}\FloatTok{0.5}
\ControlFlowTok{if}\NormalTok{ i }\KeywordTok{in}\NormalTok{ l1:}
\NormalTok{  pos}\OperatorTok{=}\NormalTok{l1.index(i) }
  \BuiltInTok{print}\NormalTok{(}\SpecialStringTok{f"}\SpecialCharTok{\{}\NormalTok{i}\SpecialCharTok{\}}\SpecialStringTok{ se encuentra en la posición }\SpecialCharTok{\{}\NormalTok{pos}\SpecialCharTok{\}}\SpecialStringTok{"}\NormalTok{)}
\ControlFlowTok{else}\NormalTok{:}
    \BuiltInTok{print}\NormalTok{(}\SpecialStringTok{f"}\SpecialCharTok{\{}\NormalTok{i}\SpecialCharTok{\}}\SpecialStringTok{ no está en la lista"}\NormalTok{)}
\end{Highlighting}
\end{Shaded}
\end{frame}

\begin{frame}[fragile]{Capturar errores con \texttt{try-except}}
\phantomsection\label{capturar-errores-con-try-except}
\begin{tcolorbox}[enhanced jigsaw, title=\textcolor{quarto-callout-tip-color}{\faLightbulb}\hspace{0.5em}{Tip}, bottomtitle=1mm, opacityback=0, bottomrule=.15mm, arc=.35mm, colframe=quarto-callout-tip-color-frame, colback=white, opacitybacktitle=0.6, colbacktitle=quarto-callout-tip-color!10!white, left=2mm, coltitle=black, leftrule=.75mm, rightrule=.15mm, toptitle=1mm, titlerule=0mm, breakable, toprule=.15mm]

Como se La forma \textbf{técnica} de tratar con este error en python es
capturándolo con \texttt{try-except}:

\begin{Shaded}
\begin{Highlighting}[]
\NormalTok{i}\OperatorTok{=}\FloatTok{0.5}
\ControlFlowTok{try}\NormalTok{:}
\NormalTok{    pos }\OperatorTok{=}\NormalTok{ l1.index(i)}
    \BuiltInTok{print}\NormalTok{(}\SpecialStringTok{f"}\SpecialCharTok{\{}\NormalTok{i}\SpecialCharTok{\}}\SpecialStringTok{ se encuentra en la posición }\SpecialCharTok{\{}\NormalTok{pos}\SpecialCharTok{\}}\SpecialStringTok{"}\NormalTok{)}
\ControlFlowTok{except} \PreprocessorTok{ValueError}\NormalTok{:}
    \BuiltInTok{print}\NormalTok{(}\SpecialStringTok{f"}\SpecialCharTok{\{}\NormalTok{i}\SpecialCharTok{\}}\SpecialStringTok{ no está en la lista"}\NormalTok{)}
\end{Highlighting}
\end{Shaded}

Con \texttt{try}, python intenta ejecutar el código dentro del bloque.
Si ocurre un \texttt{ValueError}, python salta al bloque \texttt{except}
y ejecuta ese código.

\end{tcolorbox}
\end{frame}

\section{Sub-listas}\label{sub-listas}

\begin{frame}[fragile]{Sublistas}
\phantomsection\label{sublistas}
\index{Listas!Sublistas}

Para extraer una sublista de una lista empleamos la sintaxis
\texttt{l{[}i:j{]}}, \texttt{i} corresponde al índice del elemento
inicial a partir del cual, incluido, comenzará la sublista, mientras que
\texttt{j} indica el extremo superior, la sublista llegará hasta el
elemento con índice \texttt{j-1}. La sentencia \texttt{l{[}i:{]}}
regresa una sublista de \texttt{l} que inicia en el índice \texttt{i}
hasta el final de la lista, de manera análoga \texttt{l{[}:j{]}} es una
sublista que inicia en el índice cero de la lista hasta el índice
\texttt{j-1}, por último \texttt{l{[}:{]}} regresa la lista completa.
\end{frame}

\begin{frame}[fragile]{Sublistas}
\phantomsection\label{sublistas-1}
Veamos

\begin{Shaded}
\begin{Highlighting}[]
\NormalTok{l1}
\end{Highlighting}
\end{Shaded}

\begin{verbatim}
[-6, -5, -4, -3, -2, -1, 1, 2, 3, 4, 5, 6]
\end{verbatim}

\phantomsection\label{annotated-cell-25}%
\begin{Shaded}
\begin{Highlighting}[]
\NormalTok{l1[}\DecValTok{2}\NormalTok{:}\DecValTok{7}\NormalTok{] }\hspace*{\fill}\NormalTok{\circled{1}}
\NormalTok{l1[}\DecValTok{2}\NormalTok{:] }\hspace*{\fill}\NormalTok{\circled{2}}
\NormalTok{l1[:}\DecValTok{4}\NormalTok{] }\hspace*{\fill}\NormalTok{\circled{3}}
\NormalTok{l1[:] }\hspace*{\fill}\NormalTok{\circled{4}}
\end{Highlighting}
\end{Shaded}

\begin{description}
\tightlist
\item[\circled{1}]
Desde el elemento con índice 2 hasta el elemento con indice 6.
\item[\circled{2}]
Desde el elemento con índice 2 hasta el último.
\item[\circled{3}]
Desde el primer (índice 0) hasta el elemento con índice 3.
\item[\circled{4}]
La lista completa.
\end{description}
\end{frame}

\begin{frame}[fragile]{Otras formas de crear listas}
\phantomsection\label{sec-ListasOtrasFormas}
\index{Listas!Creación de}

Es posible construir listas a partir de \texttt{tuplas},
\texttt{diccionarios}, \texttt{sets} (conjuntos) y otros objetos,
mediante la función \texttt{list()}. Veamos un ejemplo de creación de
listas a partir de tuplas

\begin{Shaded}
\begin{Highlighting}[]
\NormalTok{t1}\OperatorTok{=}\NormalTok{(}\DecValTok{3}\NormalTok{,}\DecValTok{4}\NormalTok{) }
\NormalTok{lt1}\OperatorTok{=}\BuiltInTok{list}\NormalTok{(t1) }\CommentTok{\# se crea lista a partir de la tupla t1 }
\BuiltInTok{print}\NormalTok{(}\BuiltInTok{type}\NormalTok{(lt1)) }
\end{Highlighting}
\end{Shaded}

\begin{verbatim}
<class 'list'>
\end{verbatim}
\end{frame}

\begin{frame}[fragile]{Combinación de las funciones \texttt{list()} y
\texttt{range()}.}
\phantomsection\label{combinaciuxf3n-de-las-funciones-list-y-range.}
\index{Listas!Función!\texttt{list()}}

\index{Listas!Función!\texttt{range()}}\index{Función!\texttt{range()}}

En el siguiente ejemplo se usa la función \texttt{range()} que
representa una secuencia \emph{inmutable} (como las tuplas),
estrictamente creciente o decreciente, de números. Esta función se
utiliza en bucles \texttt{for} para realizar un número específico de
iteraciones en función de un valor que se reduce o aumenta en una
constante (\texttt{step}), hasta que se alcanza un límite. La función
\texttt{range()} regresa un objeto de tipo \texttt{range}, pero se puede
convertir en una lista mediante la función \texttt{list()}.
\end{frame}

\begin{frame}[fragile]{Combinación de las funciones \texttt{list()} y
\texttt{range()}.}
\phantomsection\label{combinaciuxf3n-de-las-funciones-list-y-range.-1}
Ilustremos su uso con ejemplos

\phantomsection\label{annotated-cell-27}%
\begin{Shaded}
\begin{Highlighting}[]
\BuiltInTok{print}\NormalTok{(}\BuiltInTok{range}\NormalTok{(}\DecValTok{9}\NormalTok{),}\BuiltInTok{type}\NormalTok{(}\BuiltInTok{range}\NormalTok{(}\DecValTok{9}\NormalTok{))) }\hspace*{\fill}\NormalTok{\circled{1}}
\BuiltInTok{list}\NormalTok{(}\BuiltInTok{range}\NormalTok{(}\DecValTok{9}\NormalTok{)) }\hspace*{\fill}\NormalTok{\circled{2}}
\BuiltInTok{list}\NormalTok{(}\BuiltInTok{range}\NormalTok{(}\DecValTok{3}\NormalTok{,}\DecValTok{22}\NormalTok{,}\DecValTok{2}\NormalTok{)) }\hspace*{\fill}\NormalTok{\circled{3}}
\end{Highlighting}
\end{Shaded}

\begin{description}
\tightlist
\item[\circled{1}]
La clase del objeto \texttt{range(0,\ 9)} es
\texttt{\textquotesingle{}range\textquotesingle{}}
\item[\circled{2}]
Lista que contiene secuencia de enteros desde cero (0) hasta ocho (8).
\item[\circled{3}]
Lista que contiene secuencia de enteros desde 3 hasta 21 incrementando
en 2.
\end{description}
\end{frame}

\begin{frame}[fragile]{Combinación de las funciones \texttt{list()} y
\texttt{range()}.}
\phantomsection\label{combinaciuxf3n-de-las-funciones-list-y-range.-2}
\begin{itemize}
\item
  En \texttt{range(start,stop,step)}, \texttt{start} es el parámetro
  inicio, entero a partir del cual se devolverá la secuencia,
\item
  \texttt{stop} es el entero que indica el final; la secuencia termina
  en el mayor entero que es menor o igual que \texttt{stop-1},
\item
  Por último \texttt{step} es el paso, que determina el incremento entre
  cada entero de la secuencia.
\item
  Se puede iterar sobre el objeto (\texttt{range}) sin necesidad de
  convertirlo a una lista, uso que es bastante frecuente cuando se
  programa en python
\end{itemize}
\end{frame}

\begin{frame}[fragile]{Combinación de las funciones \texttt{list()} y
\texttt{range()}.}
\phantomsection\label{combinaciuxf3n-de-las-funciones-list-y-range.-3}
\begin{Shaded}
\begin{Highlighting}[]
\ControlFlowTok{for}\NormalTok{ i }\KeywordTok{in} \BuiltInTok{range}\NormalTok{(}\DecValTok{3}\NormalTok{,}\DecValTok{15}\NormalTok{,}\DecValTok{3}\NormalTok{):}
  \BuiltInTok{print}\NormalTok{(i)}
\end{Highlighting}
\end{Shaded}
\end{frame}

\section{Comprensión de listas}\label{comprensiuxf3n-de-listas}

\begin{frame}[fragile]{Comprensión de listas}
\phantomsection\label{sec-ComprensionDeListas}
\index{Listas!Comprensión de}

La comprensión de listas es una forma elegante de definir y crear listas
basadas en listas existentes. Es una forma concisa de generar una lista
al aplicar una operación a los elementos de otra lista, esto se ilustra
en el siguiente ejemplo:

\begin{Shaded}
\begin{Highlighting}[]
\NormalTok{[}\OperatorTok{{-}}\DecValTok{2}\OperatorTok{+}\FloatTok{0.5}\OperatorTok{*}\NormalTok{x }\ControlFlowTok{for}\NormalTok{ x }\KeywordTok{in} \BuiltInTok{range}\NormalTok{(}\DecValTok{9}\NormalTok{)]}
\end{Highlighting}
\end{Shaded}

\begin{verbatim}
[-2.0, -1.5, -1.0, -0.5, 0.0, 0.5, 1.0, 1.5, 2.0]
\end{verbatim}

Lo que se hace en el código anterior es evaluar la expresión \(-2+0.5x\)
para cada valor \(x\) en la secuencia generada por la función
\texttt{range()}.
\end{frame}

\begin{frame}[fragile]{Comprensión de listas (cont)}
\phantomsection\label{comprensiuxf3n-de-listas-cont}
Otro ejemplo, suponga que tiene una lista con valores de temperatura
expresada en grados centígrado (\(C\)) y se le pide que haga la
conversión a grados Fahrenheit (\(F\)) mediante la transformación
\(F=\frac 9 5 C+32\), podemos hacer la conversión de la siguiente forma

\begin{Shaded}
\begin{Highlighting}[]
\NormalTok{C}\OperatorTok{=}\NormalTok{[}\DecValTok{25}\NormalTok{,}\DecValTok{32}\NormalTok{,}\DecValTok{36}\NormalTok{,}\DecValTok{19}\NormalTok{,}\DecValTok{20}\NormalTok{,}\DecValTok{24}\NormalTok{] }\CommentTok{\#temperaturas en grados centígrados}
\NormalTok{F}\OperatorTok{=}\NormalTok{[}\FloatTok{9.0}\OperatorTok{/}\DecValTok{5}\OperatorTok{*}\NormalTok{x}\OperatorTok{+}\DecValTok{32} \ControlFlowTok{for}\NormalTok{ x }\KeywordTok{in}\NormalTok{ C] }\CommentTok{\#temperaturas en grados Fahrenheit}
\NormalTok{F}
\end{Highlighting}
\end{Shaded}

\begin{verbatim}
[77.0, 89.6, 96.8, 66.2, 68.0, 75.2]
\end{verbatim}
\end{frame}

\begin{frame}[fragile]{Comprensión de listas (cont)}
\phantomsection\label{comprensiuxf3n-de-listas-cont-1}
En la comprensión de listas se pueden utilizar sentencias condicionales
para generar nuevas listas, veamos un ejemplo: de la lista \texttt{l3}
obtener una sub-lista que contenga sólo los elementos positivos.

\begin{Shaded}
\begin{Highlighting}[]
\BuiltInTok{print}\NormalTok{(l3) }
\end{Highlighting}
\end{Shaded}

\begin{verbatim}
[2, 5, 4, -1, 4, 3, 0, -2]
\end{verbatim}

\begin{Shaded}
\begin{Highlighting}[]
\NormalTok{[x }\ControlFlowTok{for}\NormalTok{ x }\KeywordTok{in}\NormalTok{ l3 }\ControlFlowTok{if}\NormalTok{ x}\OperatorTok{\textgreater{}}\DecValTok{0}\NormalTok{]}
\end{Highlighting}
\end{Shaded}

\begin{verbatim}
[2, 5, 4, 4, 3]
\end{verbatim}
\end{frame}

\section{\texorpdfstring{Copiar listas vs.~asignar referencias
(\texttt{Aliasing})}{Copiar listas vs.~asignar referencias (Aliasing)}}\label{copiar-listas-vs.-asignar-referencias-aliasing}

\begin{frame}[fragile]{\texttt{listaB=listaA}}
\phantomsection\label{listablistaa}
Tenemos la lista \texttt{l3}

\begin{verbatim}
[2, 5, 4, -1, 4, 3, 0, -2]
\end{verbatim}

Igualamos \texttt{l4\ a\ l3}

\begin{Shaded}
\begin{Highlighting}[]
\NormalTok{l4}\OperatorTok{=}\NormalTok{l3  }
\end{Highlighting}
\end{Shaded}

Modificamos \texttt{l4} agregando \texttt{6} al final

\begin{Shaded}
\begin{Highlighting}[]
\NormalTok{l4.append(}\DecValTok{6}\NormalTok{)}
\end{Highlighting}
\end{Shaded}

Imprimimos \texttt{l3}

\begin{Shaded}
\begin{Highlighting}[]
\BuiltInTok{print}\NormalTok{(l3) }
\end{Highlighting}
\end{Shaded}

\begin{verbatim}
[2, 5, 4, -1, 4, 3, 0, -2, 6]
\end{verbatim}
\end{frame}

\begin{frame}[fragile]{\texttt{listaB=listaA.copy()}}
\phantomsection\label{listablistaa.copy}
Tenemos la lista \texttt{l3}

\begin{verbatim}
[2, 5, 4, -1, 4, 3, 0, -2, 6]
\end{verbatim}

Hacemos \texttt{l4=l3.copy()}

\begin{Shaded}
\begin{Highlighting}[]
\NormalTok{l4}\OperatorTok{=}\NormalTok{l3.copy() }
\end{Highlighting}
\end{Shaded}

Modificamos \texttt{l4} agregando \texttt{7} al final

\begin{Shaded}
\begin{Highlighting}[]
\NormalTok{l4.append(}\DecValTok{7}\NormalTok{)}
\end{Highlighting}
\end{Shaded}

Imprimimos \texttt{l3}

\begin{Shaded}
\begin{Highlighting}[]
\BuiltInTok{print}\NormalTok{(l3)  }
\end{Highlighting}
\end{Shaded}

\begin{verbatim}
[2, 5, 4, -1, 4, 3, 0, -2, 6]
\end{verbatim}
\end{frame}

\begin{frame}[fragile]{Importante}
\phantomsection\label{importante}
\begin{tcolorbox}[enhanced jigsaw, title=\textcolor{quarto-callout-tip-color}{\faLightbulb}\hspace{0.5em}{Tip}, bottomtitle=1mm, opacityback=0, bottomrule=.15mm, arc=.35mm, colframe=quarto-callout-tip-color-frame, colback=white, opacitybacktitle=0.6, colbacktitle=quarto-callout-tip-color!10!white, left=2mm, coltitle=black, leftrule=.75mm, rightrule=.15mm, toptitle=1mm, titlerule=0mm, breakable, toprule=.15mm]

\begin{itemize}
\item
  Si se hace \texttt{lista2\ =\ lista1}, lo que se conoce como
  \texttt{creación\ de\ alias} \texttt{(Aliasing)} o
  \texttt{asignación\ por\ referencia}, al modificar \texttt{lista2}
  también alterará \texttt{lista1}.
\item
  El método \texttt{.copy()} (ej. \texttt{lista2\ =\ lista1.copy()})
  crea \texttt{lista2} como un objeto independiente. Esto se conoce como
  \texttt{copia\ superficial} o \texttt{Shallow\ Copy}
\end{itemize}

\end{tcolorbox}
\end{frame}

\section{Ejemplo de aplicación}\label{ejemplo-de-aplicaciuxf3n}

\begin{frame}[fragile]{Aplicación en control de calidad}
\phantomsection\label{aplicaciuxf3n-en-control-de-calidad}
Supongamos que monitoreamos el grosor de piezas en un proceso de
manufactura. Almacenamos las mediciones de las muestras en una lista.

\begin{Shaded}
\begin{Highlighting}[]
\CommentTok{\# 1. Recolección de datos}
\NormalTok{mediciones }\OperatorTok{=}\NormalTok{ [}\FloatTok{10.2}\NormalTok{, }\FloatTok{10.1}\NormalTok{, }\FloatTok{9.8}\NormalTok{, }\FloatTok{10.3}\NormalTok{, }\FloatTok{11.5}\NormalTok{, }\FloatTok{9.9}\NormalTok{, }\FloatTok{10.0}\NormalTok{, }\FloatTok{10.1}\NormalTok{, }\FloatTok{8.5}\NormalTok{, }\FloatTok{10.2}\NormalTok{]}

\CommentTok{\# 2. Uso de funciones nativas para calcular la Media}
\NormalTok{n }\OperatorTok{=} \BuiltInTok{len}\NormalTok{(mediciones)}
\NormalTok{suma\_total }\OperatorTok{=} \BuiltInTok{sum}\NormalTok{(mediciones)}
\NormalTok{media\_muestral }\OperatorTok{=}\NormalTok{ suma\_total }\OperatorTok{/}\NormalTok{ n}

\BuiltInTok{print}\NormalTok{(}\SpecialStringTok{f"La media del proceso es: }\SpecialCharTok{\{}\NormalTok{media\_muestral}\SpecialCharTok{:.3f\}}\SpecialStringTok{"}\NormalTok{)}
\end{Highlighting}
\end{Shaded}

\begin{verbatim}
La media del proceso es: 10.060
\end{verbatim}
\end{frame}

\begin{frame}[fragile]{Identificando anomalías}
\phantomsection\label{identificando-anomaluxedas}
Si una medición está por encima de \(\overline x +0.5\) o por debajo de
\(\overline x - 0.5\) se considera una anomalía y debemos identificarlas

\begin{Shaded}
\begin{Highlighting}[]
\CommentTok{\# Establecemos límites de control empíricos}
\NormalTok{limite\_sup }\OperatorTok{=}\NormalTok{ media\_muestral }\OperatorTok{+} \FloatTok{0.5}
\NormalTok{limite\_inf }\OperatorTok{=}\NormalTok{ media\_muestral }\OperatorTok{{-}} \FloatTok{0.5}

\CommentTok{\# Extraemos solo las piezas que están fuera de control}
\NormalTok{piezas\_defectuosas }\OperatorTok{=}\NormalTok{ [x }\ControlFlowTok{for}\NormalTok{ x }\KeywordTok{in}\NormalTok{ mediciones }\ControlFlowTok{if}\NormalTok{ x }\OperatorTok{\textgreater{}}\NormalTok{ limite\_sup }\KeywordTok{or}\NormalTok{ x }\OperatorTok{\textless{}}\NormalTok{ limite\_inf]}

\BuiltInTok{print}\NormalTok{(}\SpecialStringTok{f"Mediciones fuera de control encontradas: }\SpecialCharTok{\{}\NormalTok{piezas\_defectuosas}\SpecialCharTok{\}}\SpecialStringTok{"}\NormalTok{)}
\end{Highlighting}
\end{Shaded}

\begin{verbatim}
Mediciones fuera de control encontradas: [11.5, 8.5]
\end{verbatim}

\begin{Shaded}
\begin{Highlighting}[]
\BuiltInTok{print}\NormalTok{(}\SpecialStringTok{f"Total de defectos: }\SpecialCharTok{\{}\BuiltInTok{len}\NormalTok{(piezas\_defectuosas)}\SpecialCharTok{\}}\SpecialStringTok{"}\NormalTok{)}
\end{Highlighting}
\end{Shaded}

\begin{verbatim}
Total de defectos: 2
\end{verbatim}
\end{frame}

\section{Listas anidadas}\label{listas-anidadas}

\begin{frame}[fragile]{Listas anidadas}
\phantomsection\label{listas-anidadas-1}
\index{Listas!Anidadas}

Como con las tuplas, es posible tener dentro de una lista otras listas o
tuplas u objetos de distinto tipo, a continuación se crea una lista que
\emph{anida} las listas \texttt{C} y \texttt{F} creadas anteriormente

\begin{Shaded}
\begin{Highlighting}[]
\NormalTok{T}\OperatorTok{=}\NormalTok{[C,F]}
\end{Highlighting}
\end{Shaded}
\end{frame}

\begin{frame}[fragile]{Listas anidadas (cont)}
\phantomsection\label{listas-anidadas-cont}
Para recuperar el tercer elemento de la lista \texttt{F}, que está
anidada en \texttt{T} se usa la indexación, recuperando primero a
\texttt{F} y luego al elemento dentro de \texttt{F}

\begin{Shaded}
\begin{Highlighting}[]
\NormalTok{T[}\DecValTok{1}\NormalTok{][}\DecValTok{2}\NormalTok{]}
\end{Highlighting}
\end{Shaded}

\begin{verbatim}
96.8
\end{verbatim}
\end{frame}

\begin{frame}[fragile]{Ordenar una lista.}
\phantomsection\label{ordenar-una-lista.}
\index{Listas!Ordenar una} \index{Método!\texttt{sort()}}

El método \texttt{sort} ordena los elementos de una lista, por defecto
lo hace de manera ascendente, es decir, de menor a mayor, si se quiere
que el orden sea descendente usamos la opción \texttt{reverse=True}.

\begin{Shaded}
\begin{Highlighting}[]
\NormalTok{l3}
\end{Highlighting}
\end{Shaded}

\begin{verbatim}
[2, 5, 4, -1, 4, 3, 0, -2, 6]
\end{verbatim}

\begin{Shaded}
\begin{Highlighting}[]
\NormalTok{l3.sort()}
\NormalTok{l3.sort(reverse}\OperatorTok{=}\VariableTok{True}\NormalTok{) }\CommentTok{\#orden descendente }
\end{Highlighting}
\end{Shaded}
\end{frame}




\end{document}
